% Volume II: Causal Explosion
% Part IV. Time, Delay, and the Maturation of Evidence
% Chapter 22: Statutes, Reopening, and Finality
% Spine: proof-spines/ch022-spine.md

\chapter{Statutes, Reopening, and Finality}
\label{ch:ch022}

Law must reconcile delayed learning with the need for closure. This chapter models that problem through the loss function
\[
L(T)
=
L_{\mathrm{premature}}(T)
+
L_{\mathrm{delay}}(T),
\]
\ToyModel\ This loss decomposition is a stylized design model for closure rather than a calibrated doctrinal metric.
where \(T\) is the time of closure. Finality reached too early locks in avoidable error; finality deferred indefinitely destroys the security and administrability that judgment requires.

The optimal \(T\) is therefore neither immediate nor infinite. It varies with the reversibility and severity of the action under review, and it explains why statutes, reopening standards, clemency, and other second-look procedures should be understood as design variables rather than mere doctrinal leftovers. A closure rule set by \(L(T)\) is distinct from a closure that never arrives because no one revisits the file: the latter is \textbf{pathological delay}, produced by indifference or dysfunction rather than by any evidentiary work being done in the interval, and it is not an instance of the restraint this chapter otherwise defends. The chapter ends the Part by treating finality as a problem of structured corrigibility: how to remain open to matured evidence without making decision impossible.
